% unidades/26-citar-reportar.tex
\chapter{📢 \ruby{引用}{いんよう}と\ruby{報告}{ほうこく} — Citar y Reportar}
\unidadheader{26}{引用・報告}{〜という・〜といわれている・〜とのことだ}{Transmitir información de otras fuentes}

\begin{tcolorbox}[vocabbox, title={\emoji{pushpin} Vocabulario}]
\begin{tabularx}{\linewidth}{llX}
\toprule
\textbf{Japonés} & \textbf{Español} & \textbf{Tipo} \\ \midrule
\vocabitem{噂}{うわさ}{rumor / habladuría}{sustantivo}
\vocabitem{伝統}{でんとう}{tradición}{sustantivo}
\vocabitem{報告する}{ほうこくする}{reportar / informar}{v. suru}
\vocabitem{確認する}{かくにんする}{confirmar / verificar}{v. suru}
\vocabitem{伝える}{つたえる}{transmitir / comunicar}{v. G2}
\vocabitem{情報}{じょうほう}{información}{sustantivo}
\vocabitem{予定}{よてい}{plan / programa}{sustantivo}
\bottomrule
\end{tabularx}
\end{tcolorbox}

\subsection*{\emoji{speech-balloon} Frases Clave}
\frase{\ruby{山田}{やまだ}さんという\ruby{人}{ひと}を\ruby{知}{し}っていますか？}{¿Conoce a una persona llamada Yamada?}
\frase{\ruby{日本}{にほん}は\ruby{安全}{あんぜん}だといわれています。}{Se dice que Japón es seguro.}
\frase{\ruby{会議}{かいぎ}は\ruby{延期}{えんき}になるとのことです。}{Según me informaron, la reunión se pospondrá.}
\frase{\ruby{彼女}{かのじょ}は\ruby{来}{こ}ないと\ruby{聞}{き}きました。}{Escuché que ella no vendrá.}

\subsection*{\emoji{gear} Gramática}
\patron{という — «llamado X» o «que dice X»}
  {N + という + N \quad / \quad [plain clause] + という}
  {\ej{「\ruby{桜}{さくら}」という\ruby{花}{はな}を\ruby{知}{し}っていますか。}{¿Conoces una flor llamada «sakura»?}
   \ej{\ruby{明日}{あした}\ruby{休}{やす}みだという\ruby{話}{はなし}を\ruby{聞}{き}いた。}{Escuché que mañana es día libre.}
   \ej{\ruby{彼}{かれ}が\ruby{合格}{ごうかく}したというニュースは\ruby{本当}{ほんとう}ですか。}{¿Es verdad la noticia de que él aprobó?}}

\patron{といわれている — «se dice que» (conocimiento general)}
  {[plain clause] + といわれている / とされている}
  {\ej{\ruby{緑茶}{りょくちゃ}は\ruby{体}{からだ}にいいといわれています。}{Se dice que el té verde es bueno para el cuerpo.}
   \ej{\ruby{日本語}{にほんご}は\ruby{難}{むずか}しい\ruby{言語}{げんご}だといわれています。}{Se dice que el japonés es un idioma difícil.}
   \ej{この\ruby{地域}{ちいき}は\ruby{昔}{むかし}からいい\ruby{場所}{ばしょ}だとされています。}{Se considera que esta zona ha sido un buen lugar desde tiempos antiguos.}}

\patron{とのことだ / と聞いている — reporte indirecto}
  {[plain clause] + とのことだ (formal) / と\ruby{聞}{き}いている / と\ruby{言}{い}っていた}
  {\ej{\ruby{社長}{しゃちょう}は\ruby{明日}{あした}いらっしゃらないとのことです。}{Según me informaron, el director no estará mañana.}
   \ej{\ruby{電車}{でんしゃ}が\ruby{遅}{おく}れていると\ruby{聞}{き}いています。}{He oído que el tren está retrasado.}
   \ej{\ruby{田中}{たなか}さんは\ruby{転勤}{てんきん}するとおっしゃっていた。}{Según dijo el sr. Tanaka, se trasladará a otra oficina.}}

\comparacion{といわれている vs らしい}
  {といわれている\\\small conocimiento general establecido\\\small \ruby{体}{からだ}にいいといわれている\\\textit{(es sabido / se dice generalmente)}}
  {らしい\\\small evidencia reciente / impresión\\\small \ruby{体}{からだ}にいいらしい\\\textit{(parece / da la impresión de que)}}
  {といわれている implica que muchas personas lo dicen o que es un hecho cultural aceptado. らしい es más personal e indica que el hablante percibe evidencia directa o ha escuchado algo recientemente. とのことだ es más formal y se usa al transmitir información de una fuente específica.}

\coloquial{〜って (contracción casual de という/と)}{
  En habla casual, という se contrae a って: 「山田さんって人、知ってる？」「明日休みって聞いた。」「難しいって言ってたよ。」 Esta contracción って es omnipresente en japonés coloquial como marcador de cita.}

\culturabox{\ruby{伝言}{でんごん}と\ruby{日本語}{にほんご} — Reportar en contexto japonés}{
  El japonés tiene un sistema rico para distinguir qué sabe el hablante de primera mano (と思う), qué le contaron (と聞いた、とのことだ) y qué es conocimiento cultural establecido (といわれている). Esta distinción epistémica — quién sabe qué y de dónde — es fundamental en la comunicación japonesa formal, especialmente en contextos de negocios y academia.
}

\lectura{\ruby{うわさ話}{うわさばなし}\\[0.4em]
  「\ruby{田中}{たなか}さんって\ruby{来月}{らいげつ}\ruby{結婚}{けっこん}するんだって！」「え、\ruby{本当}{ほんとう}に？」「うん、\ruby{友達}{ともだち}から\ruby{聞}{き}いたんだけど。」「でも\ruby{噂}{うわさ}でしょ？\ruby{確認}{かくにん}したほうがいいと\ruby{思}{おも}う。」「そうだね。\ruby{直接}{ちょくせつ}\ruby{聞}{き}いてみる。」}
{Rumores\\[0.4em]
  «¡Dicen que el sr. Tanaka se va a casar el mes que viene!» «¿De verdad?» «Sí, lo escuché de un amigo.» «Pero es un rumor, ¿no? Creo que es mejor confirmarlo.» «Tienes razón. Lo preguntaré directamente.»}

\begin{tcolorbox}[ejerciciobox, title={\emoji{pencil} Ejercicios}]
\begin{enumerate}
  \item Usa という: presenta la palabra 「\ruby{侘}{わ}び\ruby{寂}{さ}び」 como concepto japonés.
  \item Convierte a といわれている: «El café es bueno para la concentración.»
  \item ¿Cuándo usas とのことだ vs と\ruby{聞}{き}いた? Describe la diferencia.
\end{enumerate}
\textbf{Respuestas:}
1) 「\ruby{侘}{わ}び\ruby{寂}{さ}び」という\ruby{日本}{にほん}の\ruby{概念}{がいねん}があります \quad
2) コーヒーは\ruby{集中力}{しゅうちゅうりょく}にいいといわれています \quad
3) とのことだ: \ruby{上司}{じょうし}など\ruby{特定}{とくてい}の\ruby{人}{ひと}からの\ruby{伝達}{でんたつ} (formal); と\ruby{聞}{き}いた: \ruby{一般的}{いっぱんてき}な\ruby{情報}{じょうほう} (casual)
\end{tcolorbox}

\begin{tcolorbox}[kanjibox, title={\emoji{mahjong} Kanjis}]
\begin{tabularx}{\linewidth}{cllXX}
\toprule
\textbf{Kanji} & \textbf{On} & \textbf{Kun} & \textbf{Significado} & \textbf{Vocab} \\ \midrule
\kanjirow{噂}{そん}{\ruby{うわさ}{}}{rumor}{\ruby{噂}{うわさ} / \ruby{噂話}{うわさばなし}}
\kanjirow{伝}{でん}{\ruby{つた}{える}}{transmitir}{\ruby{伝}{つた}える / \ruby{伝統}{でんとう}}
\kanjirow{報}{ほう}{\ruby{むく}{いる}}{reportar}{\ruby{報告}{ほうこく} / \ruby{情報}{じょうほう}}
\kanjirow{確}{かく}{—}{confirmar}{\ruby{確認}{かくにん} / \ruby{確実}{かくじつ}}
\bottomrule
\end{tabularx}
\end{tcolorbox}
